\begin{figure}[H]
\centering
\begin{minipage}[t]{0.45\textwidth}
  \vspace{0pt}\centering
  \includegraphics[width=\linewidth]{figures/ch05/panels/fig5_1a_contraction.pdf}\\[-0.2em]
  {\small $\displaystyle \frac{dX_t}{dt}=-\alpha X_t$}
\end{minipage}\hfill
\begin{minipage}[t]{0.45\textwidth}
  \vspace{0pt}\centering
  \includegraphics[width=\linewidth]{figures/ch05/panels/fig5_1b_brownian.pdf}\\[-0.2em]
  {\small $\displaystyle W_{t+\Delta t}-W_t\sim\mathcal N(0,\Delta t)$}
\end{minipage}

%\vspace{0.8em}

\begin{minipage}[t]{0.45\textwidth}
  \vspace{0pt}\centering
  \includegraphics[width=\linewidth]{figures/ch05/panels/fig5_1c_ou.pdf}\\[-0.2em]
  {\small $\displaystyle dX_t=-\alpha X_t\,dt+\sigma\,dW_t$}
\end{minipage}\hfill
\begin{minipage}[t]{0.45\textwidth}
  \vspace{0pt}
  \caption{\textbf{三种基本的连续时间运动规律。}
  确定性收缩产生平滑轨迹并逐渐衰减到原点。Brownian 运动具有零均值的独立高斯增量，因此轨迹随机游走，其典型扩散尺度与 $\sqrt{t}$ 成正比；虚线表示单位 Brownian 运动的 $\pm2\sqrt{t}$。Ornstein--Uhlenbeck 过程把随机波动与指向原点的回复漂移结合起来；虚线表示平稳标准差的两倍。}
  \label{fig:continuous-three-motion-laws}
\end{minipage}
\end{figure}
